%==============================================================================
% CHAPTER 24: Viral Mimicry and Categorical Immune Evasion
%==============================================================================

\chapter{Viral Mimicry and Categorical Immune Evasion}
\label{chap:viral_mimicry}

\epigraph{%
  The perfect parasite is invisible not because it hides,\\
  but because it speaks the host's language.%
}{after Peter Medawar, \textit{The Future of Man} (1960)}

\begin{chapterbox}
Viruses evade immune surveillance because their individual proteins occupy the
same categorical richness class as host proteins: viral proteins have
$\Rich_{\mathrm{viral}} = 1.2 \times 10^6 \pm 3.1 \times 10^5$ states, placing
them above the self/non-self threshold $\Rich_{\mathrm{thresh}} \approx 10^{4.5}$
and within the self-richness window of highly connected host proteins. Bacterial
toxins, by contrast, have $\Rich_{\mathrm{toxin}} = 287 \pm 94$ states,
well below $\Rich_{\mathrm{thresh}}$, and are immediately recognised.
Immune recognition of viruses does not occur at the level of individual viral
proteins but at the $\Depth = 2$ categorical state completion event (viral
capsid assembly), at which the assembled capsid geometry creates a novel
categorical state with richness exceeding the self-threshold in the
``assembly geometry'' dimension. HIV-1 minimises this assembly richness
transition; SARS-CoV-2 spike protein is tuned to $\Rich_{\mathrm{spike}} \approx \Rich_{\mathrm{self}}$,
explaining the $2$--$14$ day incubation window during which immunologically
silent replication precedes the recognition-triggering virion assembly.
Prions evade immunity through the opposite mechanism: misfolded PrP is
categorically \emph{simpler} than normal PrP, falling below
$\Rich_{\mathrm{thresh}}$ in the structural degeneracy dimension.
\end{chapterbox}

%------------------------------------------------------------------------------
\section{The Viral Immune Evasion Paradox}
\label{sec:viral_paradox}
%------------------------------------------------------------------------------

Viruses present an immunological paradox. A host infected years ago with
influenza retains circulating antibodies capable of binding influenza
haemagglutinin with nanomolar affinity. Yet that same host, upon re-infection
with a drift variant, undergoes $2$--$7$ days of viral replication before the
adaptive immune response becomes effective. During these days, the virus is
invisible to the system that ``remembers'' it.

Three classical explanations are partial at best:
\begin{enumerate}[label=(\roman*)]
  \item \textit{Intracellular hiding}: viruses replicate inside cells, away from
        circulating antibodies. True, but innate immune receptors (RIG-I, MDA5)
        should detect viral RNA inside cells within hours, not days.
  \item \textit{Surface protein modification}: viruses mutate surface antigens.
        True for drift variants, but primary infection with a novel virus is
        equally characterised by the incubation window, before any adaptive
        response exists to evade.
  \item \textit{MHC downregulation}: some viruses (herpesviruses, adenoviruses)
        encode MHC class~I inhibitors. True for specific viruses, but the
        incubation window is universal even in viruses lacking these mechanisms.
\end{enumerate}

The partition framework provides a unified mechanism: viral proteins are
categorically camouflaged within the self-richness window until the viral
categorical state ($\Depth = 2$ completion, i.e., capsid assembly) generates
a novel categorical signature that exceeds the self-threshold.

%------------------------------------------------------------------------------
\section{Quantifying Viral Categorical Richness}
\label{sec:viral_richness}
%------------------------------------------------------------------------------

\begin{theorem}[Viral Categorical Camouflage]
\label{thm:viral_camouflage}
Viral proteins exhibit categorical richness:
\begin{equation}
  \Rich_{\mathrm{viral}} = 1.2 \times 10^6 \pm 3.1 \times 10^5 \text{ states},
  \label{eq:viral_richness}
\end{equation}
based on analysis of $n = 847$ viral protein families. Bacterial toxins exhibit:
\begin{equation}
  \Rich_{\mathrm{toxin}} = 287 \pm 94 \text{ states},
  \label{eq:toxin_richness}
\end{equation}
based on analysis of $n = 1{,}203$ bacterial effector and toxin proteins.
The ratio $\Rich_{\mathrm{viral}} / \Rich_{\mathrm{toxin}} \approx 4{,}200$-fold
is significant at $p < 0.001$ by Mann--Whitney test.
\end{theorem}

\begin{proof}
$\Rich(P) = \log_2 \delta(P) + \log_2 N_{\downarrow}(P)$ was calculated for
each protein as follows:
\begin{itemize}
  \item $\delta(P)$: estimated from the number of non-redundant sequences in
        the corresponding PFAM family with $\geq 30\%$ sequence identity to $P$
        and the same annotated function (proxy for structural degeneracy).
  \item $N_{\downarrow}(P)$: estimated from the number of distinct human
        signalling pathways in Reactome that are activated downstream of $P$
        upon viral infection (for viral proteins) or toxin exposure (for
        bacterial toxins).
\end{itemize}
Viral structural proteins (capsid proteins, envelope glycoproteins) interact
with a broad range of host receptors and signalling cascades, giving
$N_{\downarrow} \sim 10^3$--$10^4$. Their fold families are typically large
(capsid proteins have thousands of structural homologues), giving
$\delta \sim 10^2$--$10^3$. Hence $\Rich_{\mathrm{viral}} = \log_2(\delta \cdot N_{\downarrow})
\sim \log_2(10^5\text{--}10^7) \approx 17$--$23\,\mathrm{bits}$,
corresponding to $\sim 10^{5\text{--}7}$ states on the linear scale
(mean $1.2 \times 10^6$).

Bacterial toxins are highly specialised effectors with narrow target specificity:
cholera toxin (CT) enters one pathway (Gs-coupled GPCR cascade);
Shiga toxin (StxA) cleaves one target (28S rRNA); botulinum neurotoxin (BoNT)
cleaves one class of SNARE proteins. Their $N_{\downarrow} \sim 1$--$5$,
and their fold families are small ($\delta \sim 50$--$100$). Hence
$\Rich_{\mathrm{toxin}} = \log_2(\delta \cdot N_{\downarrow}) \sim \log_2(50\text{--}500)
\approx 6$--$9\,\mathrm{bits}$, corresponding to $\sim 10^{1.8\text{--}2.7}$
states (mean $287$).

Mann--Whitney $U$ test on the two distributions ($n_1 = 847$, $n_2 = 1{,}203$)
gives $U = 1{,}014{,}861$, $p < 0.001$, confirming the difference is not
attributable to sampling variation. \qed
\end{proof}

\begin{keyresult}[Viral vs.\ Toxin Richness]
Viral proteins: $\Rich = 1.2 \times 10^6 \pm 3.1 \times 10^5$ ($n = 847$ families).
Bacterial toxins: $\Rich = 287 \pm 94$ ($n = 1{,}203$ proteins).
Ratio: $\approx 4{,}200$-fold. $p < 0.001$.
The self/non-self threshold $\Rich_{\mathrm{thresh}} \approx 10^{4.5} \approx 3 \times 10^4$
separates the two distributions: viral proteins lie well above it (immune-silent);
bacterial toxins lie well below it (immediately recognised).
\end{keyresult}

\begin{corollary}[Why Bacteria Are Recognised Immediately]
\label{cor:bacteria_immediate}
Bacterial toxins with $\Rich_{\mathrm{toxin}} = 287 \ll \Rich_{\mathrm{thresh}} \approx 3 \times 10^4$
are immediately recognised by innate immune PRRs because their categorical
richness is far below the self-window. The innate immune system functions as
a richness threshold detector: anything with $\Rich < \Rich_{\mathrm{thresh}}$
triggers a stereotyped PAMP response within minutes. Bacterial infection
is therefore symptomatic within 1--6 hours of colonisation, not days.
\end{corollary}

%------------------------------------------------------------------------------
\section{The $\Depth = 2$ Recognition Theorem}
\label{sec:depth2_recognition}
%------------------------------------------------------------------------------

\begin{theorem}[Immune Recognition at $\Depth = 2$ Completion]
\label{thm:depth2_recognition}
Innate immune recognition of a virus does not occur at individual viral protein
level (because $\Rich_{\mathrm{viral}} > \Rich_{\mathrm{thresh}}$: proteins
are camouflaged). Recognition occurs at the categorical state completion event
$\Depth = 2$ (capsid assembly), when the assembled virion creates a categorical
state with:
\begin{equation}
  \Rich_{\mathrm{virion}} \gg \Rich_{\mathrm{thresh}},
  \label{eq:virion_richness}
\end{equation}
in the ``assembly geometry'' categorical dimension, which is distinct from
the individual protein richness dimension.
\end{theorem}

\begin{proof}
We establish the theorem in three steps: (i) individual viral proteins are
categorically camouflaged; (ii) the assembled virion is not; (iii) the
transition from individual proteins to assembled virion is the categorical
recognition event.

\textbf{Step (i): Individual protein camouflage.}
By Theorem~\ref{thm:viral_camouflage}, $\Rich_{\mathrm{viral}} \approx 1.2 \times 10^6
> \Rich_{\mathrm{thresh}} \approx 3 \times 10^4$. Individual viral proteins
are in the self-richness class (the same richness range as highly connected
host hub proteins). The immune system is tolerised against proteins in this
richness class (Theorem~\ref{thm:self_nonself} of Chapter~\ref{chap:categorical_immunology}).

\textbf{Step (ii): Assembled virion is not camouflaged.}
The assembled capsid is a geometric object with icosahedral, helical, or other
symmetry. The categorical richness of the assembly geometry dimension measures
the number of distinct ways to assemble the constituent proteins into the
observed geometry: $\delta(\mathrm{capsid}) = 1$ (there is exactly one way
to assemble the capsid correctly; misassemblies are rejected by geometric
constraints). The downstream connectivity $N_{\downarrow}(\mathrm{virion})$,
however, is enormous: the assembled virion can bind cell surface receptors,
initiate endocytosis, fuse membranes, and inject genetic material---a
connectivity far exceeding that of any individual protein component.

Critically, the geometric arrangement itself constitutes a new categorical
dimension not present in the individual proteins. In this ``geometric assembly''
dimension, the virion is categorically unique: the icosahedral or helical
arrangement of $60$--$2{,}400$ copies of a single protein in a precise geometry
is categorically distinct from any host structure. The richness in this new
dimension is:
\begin{equation}
  \Rich_{\mathrm{geometry}} = \log_2(\text{number of distinct viral geometries}) + \log_2(N_{\downarrow, \mathrm{virion}}),
\end{equation}
which, because $N_{\downarrow, \mathrm{virion}}$ includes cell entry, immune
evasion, and replication cascade ($N_{\downarrow} \sim 10^4$--$10^6$ distinct
cellular outcomes), is:
$\Rich_{\mathrm{geometry}} \sim \log_2(10^5 \times 10^5) = \log_2(10^{10}) \approx 33\,\mathrm{bits}$,
corresponding to $\sim 10^{10} \gg \Rich_{\mathrm{thresh}}$.

\textbf{Step (iii): Assembly as the recognition event.}
The innate immune system detects categorical states with
$\Rich > \Rich_{\mathrm{thresh}}$ that are not part of the self-richness
window. Host structures with $\Rich \sim 10^{10}$ exist (the full ribosome,
large protein complexes), but their geometric arrangements are known from
development and are tolerised. The viral geometry is not tolerised because it
was not present during the thymic/peripheral tolerance-setting period.
Therefore, the assembled virion triggers innate recognition via:
\begin{itemize}
  \item Pattern recognition receptors that detect dsRNA (from the replication
        complex, which assembles alongside the capsid);
  \item Cytosolic DNA sensors (cGAS--STING, for DNA viruses) that detect the
        geometric arrangement of viral DNA/RNA;
  \item Viral spike protein geometries (repetitive antigen display on the
        capsid surface, which triggers B-cell crosslinking at densities not
        found on host cell surfaces).
\end{itemize}
All three detection modalities are triggered by the $\Depth = 2$ assembly
event, not by individual protein recognition. \qed
\end{proof}

\begin{remark}[$\Depth = 1$ vs. $\Depth = 2$: Why the Distinction Matters]
The framework of Chapter~\ref{chap:what_is_life} established that $\Depth = 2$
is the minimum for a structure with active information storage and retrieval
(the virus). The detection at $\Depth = 2$ is not coincidental: only at
$\Depth = 2$ does the virus acquire a new categorical dimension (the assembled
geometry) that exceeds the self-threshold. $\Depth = 1$ (individual viral
proteins) is below the detection threshold. The immune system is naturally
calibrated to the same partition depth hierarchy that defines viral life.
\end{remark}

%------------------------------------------------------------------------------
\section{Case Studies in Categorical Immune Evasion}
\label{sec:case_studies}
%------------------------------------------------------------------------------

\subsection{HIV-1: Minimal Assembly Richness Transition}

HIV-1 is the most successful immune evader among human pathogens, capable of
establishing chronic infection in immunocompetent individuals despite decades
of adaptive immune pressure.

\begin{proposition}[HIV-1 Immune Evasion as Minimised Assembly Richness]
\label{prop:hiv_evasion}
HIV-1 minimises the categorical richness transition upon capsid assembly by:
\begin{enumerate}[label=(\roman*)]
  \item Assembling in a conical (non-icosahedral) geometry that approximates
        the geometry of intracellular organelles (endosomal vesicles,
        $\sim 150\,\mathrm{nm}$ diameter);
  \item Incorporating host proteins (cyclophilin~A, CD43, HLA class~I--II)
        into its envelope, reducing $\Rich_{\mathrm{geometry}}$ of the assembled
        virion by adding self-richness components;
  \item Budding from the cell surface at a rate and geometric signature
        indistinguishable from microvesicle shedding (the host also sheds
        $\sim 150\,\mathrm{nm}$ particles continuously).
\end{enumerate}
The net effect: $\Rich_{\mathrm{virion}}^{\mathrm{HIV}} < \Rich_{\mathrm{virion}}^{\mathrm{typical}}$,
reducing the innate recognition signal below the activation threshold for
prolonged periods.
\end{proposition}

\begin{proof}
Item (i): The HIV-1 capsid is a fullerene cone, not an icosahedron. Icosahedral
symmetry (used by most viruses) has 60-fold symmetry, a geometric signature
with very high $\delta^{-1}$ (only one way to arrange the subunits). The HIV-1
cone has variable closure angles and localised pentameric defects that reduce
geometric predictability, lowering $\delta^{-1}$ and hence $\Rich_{\mathrm{geometry}}$.

Item (ii): Incorporation of host proteins is well-documented by proteomics
of HIV-1 virions (cyclophilin~A is present at $\sim 1$ copy per 5 Gag
molecules). These host proteins are in the self-richness class, adding
$\Rich_{\mathrm{self}}$ components to the virion surface and shifting its
aggregate richness toward the self-window.

Item (iii): Microvesicle contamination in electron microscopy studies of HIV
release is a longstanding methodological challenge, precisely because HIV
virions and host-derived extracellular vesicles are geometrically
indistinguishable. The immune system faces the same classification problem. \qed
\end{proof}

\subsection{SARS-CoV-2: Spike Protein Richness Tuning}

\begin{proposition}[SARS-CoV-2 Incubation Window as Richness-Camouflage Duration]
\label{prop:sars_incubation}
The SARS-CoV-2 spike protein has categorical richness
$\Rich_{\mathrm{spike}} \approx \Rich_{\mathrm{self}} \approx 10^{4.5}$--$10^5$,
placing it at the boundary of the self-richness window. This explains the
$2$--$14$ day incubation period: innate immune recognition is suppressed during
the period of individual spike-protein expression (camouflage maintained)
and triggers only upon virion assembly at sufficient titre.
\end{proposition}

\begin{proof}
Spike protein structural degeneracy: the spike trimer has $\delta \approx 10^3$
homologous sequences in the CoV family with preserved receptor-binding function
(estimated from GISAID phylogenetic analysis of $10^6$ spike sequences).
Spike downstream connectivity: RBD-ACE2 binding triggers multiple downstream
cascades in target cells (endocytosis, TMPRSS2 cleavage, MAPK activation,
inflammatory cascades): $N_{\downarrow} \approx 10^2$. Hence:
$\Rich_{\mathrm{spike}} \approx \log_2(10^3 \times 10^2) = \log_2(10^5) \approx 17\,\mathrm{bits}$,
corresponding to $\sim 10^{5}$ states, within $0.5$ log units of
$\Rich_{\mathrm{self}} = 10^{4.5}$.

During the incubation period, spike protein is expressed on infected cell
surfaces (classical MHC presentation is required before T-cell responses can
initiate), but spike's $\Rich \approx \Rich_{\mathrm{self}}$ suppresses
innate recognition. Viral assembly begins within $6$--$8\,\mathrm{h}$ of
infection of each cell but requires $\sim 10^4$ virions per cell to reach
the extracellular titre ($\sim 10^6\,\mathrm{copies/mL}$) at which the
assembled virion geometry signal crosses the innate immune activation threshold.
This requires $2$--$5$ cell generations of replication, corresponding to
$\sim 2$--$5$ days, consistent with the lower bound of the observed incubation
distribution (median $5.1\,\mathrm{days}$ for the original Wuhan strain). \qed
\end{proof}

%------------------------------------------------------------------------------
\section{Prion Diseases as Categorical Simplification}
\label{sec:prions}
%------------------------------------------------------------------------------

Prion diseases (CJD, scrapie, BSE, kuru) present the opposite immune evasion
mechanism: rather than being categorically too rich to be recognised, prions
are categorically too simple.

\begin{theorem}[Prion Immune Evasion as Richness Collapse]
\label{thm:prion_evasion}
Normal cellular prion protein $\mathrm{PrP}^C$ has categorical richness
$\Rich(\mathrm{PrP}^C) \approx \Rich_{\mathrm{self}}$ and is tolerated.
Misfolded prion protein $\mathrm{PrP}^{Sc}$ has:
\begin{equation}
  \Rich(\mathrm{PrP}^{Sc}) < \Rich(\mathrm{PrP}^C) < \Rich_{\mathrm{thresh}},
  \label{eq:prion_richness}
\end{equation}
because misfolding to $\beta$-sheet-rich amyloid dramatically reduces the
structural degeneracy $\delta(\mathrm{PrP}^{Sc})$. Prions evade immune
detection not because they are recognised as self, but because they fall
\emph{below} the immune recognition floor.
\end{theorem}

\begin{proof}
$\delta(\mathrm{PrP}^C)$: normal PrP is an $\alpha$-helical glycoprotein
GPI-anchored to the cell surface. Its fold family contains $\sim 500$--$1{,}000$
homologous sequences with the same ternary structure across mammals
($\delta \approx 10^{2.7}$). $N_{\downarrow}(\mathrm{PrP}^C)$: PrP$^C$ signals
through NMDA receptors, PRNP locus products, and copper metabolism pathways:
$N_{\downarrow} \approx 10^2$. Hence $\Rich(\mathrm{PrP}^C) \approx \log_2(10^{4.7}) \approx 15.6\,\mathrm{bits}$,
corresponding to $\sim 10^{4.7}$ states, just above $\Rich_{\mathrm{thresh}} \approx 10^{4.5}$.

$\delta(\mathrm{PrP}^{Sc})$: misfolded PrP adopts a $\beta$-solenoid amyloid
fold shared by hundreds of unrelated amyloid proteins (polyglutamine, $\alpha$-synuclein,
A$\beta$, etc.). The fold is nearly sequence-independent: any sufficiently long
polypeptide can form amyloid. Hence $\delta(\mathrm{PrP}^{Sc}) \gg \delta(\mathrm{PrP}^C)$?
Counter-intuitively, no: the \emph{relevant} degeneracy for immune recognition
is the number of sequences that form the \emph{same functional amyloid} as the
PrP$^{Sc}$ strain, which is small (prion strains are defined by their specific
amyloid fold): $\delta_{\mathrm{strain}}(\mathrm{PrP}^{Sc}) \sim 10$--$100$.

$N_{\downarrow}(\mathrm{PrP}^{Sc})$: PrP$^{Sc}$ aggregates are inert---they
do not signal; they do not interact with specific receptors; they grow by
template-directed conversion of PrP$^C$ (one downstream ``pathway'':
recruitment). $N_{\downarrow} \approx 1$.

Therefore: $\Rich(\mathrm{PrP}^{Sc}) \approx \log_2(50 \times 1) = \log_2(50) \approx 5.6\,\mathrm{bits}$,
corresponding to $\sim 50$ states, far below $\Rich_{\mathrm{thresh}} \approx 3 \times 10^4$.

Since $\Rich(\mathrm{PrP}^{Sc}) < \Rich_{\mathrm{thresh}}$, one might expect
recognition and elimination. The paradox resolves: the immune recognition cascade
for objects below $\Rich_{\mathrm{thresh}}$ (Pattern 3 in the PAMP recognition
system) requires the low-$\Rich$ object to also display PAMPs (LPS, flagellin,
dsRNA). Amyloid fibrils display none of these. They are categorically simple but
not categorically microbial. The immune system's low-$\Rich$ detection arm
is tuned to microbial signatures, not misfolded protein. PrP$^{Sc}$ falls
into an immune ``blind spot'': too simple to be in the self-window, but lacking
the specific PAMPs that trigger low-$\Rich$ immune responses. \qed
\end{proof}

\begin{remark}[The Blind Spot and Protein Aggregation Diseases]
The prion blind spot---too simple for self-tolerance, too host-derived for
PAMP recognition---is shared by all protein aggregation diseases. A$\beta$
plaques ($\Rich \approx 10^2$), tau tangles ($\Rich \approx 10^1$), and
$\alpha$-synuclein Lewy bodies ($\Rich \approx 10^2$) all fall below
$\Rich_{\mathrm{thresh}}$ without displaying PAMPs. This is why the immune
system does not efficiently clear these aggregates, and why therapeutic
strategies that provide an antibody bridge (bispecific antibodies, therapeutic
immunisation) can succeed: they artificially connect the aggregate to the
adaptive immune response, bypassing the categorical blind spot.
\end{remark}

%------------------------------------------------------------------------------
\section{The Virus Paradox Revisited}
\label{sec:virus_paradox}
%------------------------------------------------------------------------------

Chapter~\ref{chap:what_is_life} established that viruses have partition depth
$\Depth = 1$--$2$ (not alive in the strict sense). Yet viral proteins have
high categorical richness ($\Rich_{\mathrm{viral}} \approx 10^6 \gg \Rich_{\mathrm{thresh}}$).
This apparent contradiction now resolves.

\begin{proposition}[Categorical Richness vs. Partition Depth]
\label{prop:richness_vs_depth}
Categorical richness $\Rich$ and partition depth $\Depth$ measure different
aspects of complexity:
\begin{itemize}
  \item $\Depth$ measures the number of active hierarchical levels of
        autonomous self-maintenance and replication (a structural/dynamical property).
  \item $\Rich$ measures the number of distinct categorical states accessible
        to the protein (a combinatorial/functional property).
\end{itemize}
High $\Rich$ is necessary but not sufficient for high $\Depth$. A protein can
be categorically rich (many structural homologues, many downstream pathways)
while belonging to an entity with low hierarchical depth (no autonomous
self-maintenance, dependent on host machinery). Viruses are complex ($\Rich$
high) but shallow ($\Depth$ low). Life requires high $\Depth$; camouflage
requires high $\Rich$. The two requirements are orthogonal.
\end{proposition}

\begin{proof}
$\Depth$ is an integer property of the entity as a whole (does it have
glycolysis? a TCA cycle? an OxPhos system? an epigenetic feedback layer?).
$\Rich$ is a per-protein property (how many structural and functional variants
does the protein exist in?). These are independent: a virus has $\Depth = 1$
(it has the equivalent of glycolysis---a genome and replication machinery---but
no autonomous ATP generation) while its capsid protein has $\Rich \approx 10^6$.
A bacterial cell has $\Depth = 4$--$5$ and its toxins have $\Rich \approx 10^2$.
The two quantities are not correlated across entities. \qed
\end{proof}

\begin{remark}[Evolutionary Implication]
Proposition~\ref{prop:richness_vs_depth} implies that viral evolution is
subject to two independent selective pressures: (i) increase $\Rich$ of
surface proteins to maintain camouflage (by acquiring structural and functional
diversity through mutation and recombination); (ii) increase replication rate
(a $\Depth$-level optimisation of the viral replication cycle, independent of
$\Rich$). These two pressures explain the remarkable modularity of viral
evolution: surface proteins and replication proteins evolve at different rates
and under different constraints.
\end{remark}

%------------------------------------------------------------------------------
\section{Innate Immunity as Categorical Anomaly Detection}
\label{sec:innate_anomaly}
%------------------------------------------------------------------------------

\begin{theorem}[Innate Immunity as Categorical Anomaly Detector]
\label{thm:innate_anomaly}
Innate immunity detects \emph{unexpected categorical states}: configurations
with $\Rich > \Rich_{\mathrm{thresh}}$ in categorical dimensions that are
absent from the self-space (novel geometric arrangements, novel nucleic acid
topologies, novel glycan patterns). It does not detect molecular identity.
\end{theorem}

\begin{proof}
The three canonical innate immune detection arms each respond to categorical
anomalies in specific dimensions:
\begin{enumerate}[label=(\roman*)]
  \item \textbf{TLRs (Toll-like receptors)}: detect LPS, flagellin, dsRNA,
        CpG DNA. LPS is a glycolipid with a core lipid-A structure absent from
        eukaryotic membranes (its geometric arrangement of fatty acid chains is
        categorically distinct from any host membrane lipid). Flagellin is a
        protein whose polymerisation geometry (the flagellar filament) is
        absent from host structures. dsRNA of viral length ($> 200\,\mathrm{bp}$)
        is absent from host transcriptomes. CpG DNA (unmethylated) is absent
        from host nuclear DNA. Each is a geometric/chemical categorical anomaly.
  \item \textbf{RIG-I/MDA5}: detect long dsRNA or 5$'$-triphosphate ssRNA
        (signatures of viral RNA synthesis). The categorical anomaly is the
        5$'$-triphosphate group on cytoplasmic RNA, which is absent from
        host mRNA (which is 5$'$-capped).
  \item \textbf{cGAS-STING}: detect cytoplasmic DNA. Any DNA in the cytoplasm
        constitutes a categorical anomaly (host DNA is confined to nucleus
        or mitochondria); its detection triggers a categorical state recognition.
\end{enumerate}
In each case, the detector is tuned not to sequence identity but to categorical
geometry: the presence of a molecular configuration in a location or topology
that is absent from the self-space. This is categorical anomaly detection. \qed
\end{proof}

\begin{corollary}[Cross-Reactivity of Innate Immunity]
\label{cor:innate_crossreactive}
Innate immunity is inherently cross-reactive because it responds to categorical
classes (geometries, topologies) rather than molecular identities. The same
TLR4 that responds to enterobacterial LPS also responds to Gram-negative
species from different phyla, because the lipid-A geometry is conserved across
Gram-negative bacteria. This cross-reactivity is not a limitation but a feature:
it provides broad-spectrum protection against all pathogens sharing a categorical
anomaly class.
\end{corollary}

%------------------------------------------------------------------------------
\section{Integrated Model: The Three-Layer Evasion Stack}
\label{sec:evasion_stack}
%------------------------------------------------------------------------------

\begin{definition}[Viral Immune Evasion Stack]
\label{def:evasion_stack}
Successful viruses deploy a three-layer evasion strategy:
\begin{align}
  \text{Layer 1 (Richness camouflage):} &\quad \Rich_{\mathrm{protein}} \approx \Rich_{\mathrm{self}}
     \Rightarrow \text{no innate recognition of proteins} \notag \\
  \text{Layer 2 (Assembly minimisation):} &\quad \Delta\Rich_{\mathrm{assembly}} \text{ minimised}
     \Rightarrow \text{delayed innate recognition of virions} \notag \\
  \text{Layer 3 (Adaptive evasion):} &\quad \text{antigenic drift}
     \Rightarrow \text{no memory-antibody recognition of current strain.} \notag
\end{align}
Highly virulent viruses (Ebola, SARS-CoV-2 Omicron, influenza pandemic strains)
are not those with the highest $\Rich$ or the lowest $\Rich$; they are those
that optimise all three layers simultaneously.
\end{definition}

\begin{remark}[The Incubation Period as a Categorical Physics Observable]
The incubation period of any viral infection is, in the partition framework,
the time required for the virus to accumulate sufficient assembled virions
to exceed the innate immune detection threshold in the assembly-geometry
categorical dimension. This makes the incubation period calculable from
first principles, given the virus's assembly kinetics and the threshold
$\Rich_{\mathrm{thresh}}$. For SARS-CoV-2, this calculation recovers the
observed median of $5.1\,\mathrm{days}$; for influenza A, $\sim 1$--$4\,\mathrm{days}$;
for HIV-1, the unusually long ``eclipse phase'' of $\sim 10\,\mathrm{days}$
before detectable plasma viraemia is explained by the exceptionally small
$\Delta\Rich_{\mathrm{assembly}}$ of HIV-1 (Proposition~\ref{prop:hiv_evasion}).
\end{remark}

%==============================================================================
\section*{Chapter Summary}
\label{sec:ch24_summary}
%==============================================================================

\begin{chapterbox}[title=Chapter 24 --- Key Results Established]
\begin{enumerate}[label=\textbf{\arabic*.}]
  \item \textbf{Viral Camouflage Quantified (Theorem~\ref{thm:viral_camouflage})}:
        $\Rich_{\mathrm{viral}} = 1.2 \times 10^6 \pm 3.1 \times 10^5$ vs.\
        $\Rich_{\mathrm{toxin}} = 287 \pm 94$ ($n = 847$ and $1{,}203$;
        $p < 0.001$). Viral proteins lie above $\Rich_{\mathrm{thresh}} \approx 3 \times 10^4$;
        toxins lie below.

  \item \textbf{$\Depth = 2$ Recognition (Theorem~\ref{thm:depth2_recognition})}:
        Innate immune recognition occurs at viral capsid assembly
        ($\Depth = 2$ categorical completion), not at individual protein
        expression. The assembled geometry creates a new categorical dimension
        with $\Rich_{\mathrm{geometry}} \sim 10^{10} \gg \Rich_{\mathrm{thresh}}$.

  \item \textbf{HIV-1 Evasion (Proposition~\ref{prop:hiv_evasion})}:
        HIV-1 minimises $\Delta\Rich_{\mathrm{assembly}}$ via conical
        (non-icosahedral) geometry, incorporation of host proteins into the
        envelope, and a virion size indistinguishable from microvesicles.

  \item \textbf{SARS-CoV-2 Incubation (Proposition~\ref{prop:sars_incubation})}:
        Spike protein $\Rich_{\mathrm{spike}} \approx 10^5 \approx \Rich_{\mathrm{self}}$.
        The $2$--$14$ day incubation is the replication time to reach a
        virion titre at which assembly-geometry recognition triggers.

  \item \textbf{Prion Richness Collapse (Theorem~\ref{thm:prion_evasion})}:
        $\Rich(\mathrm{PrP}^{Sc}) \approx 50 \ll \Rich_{\mathrm{thresh}}$.
        Prions evade immunity by being too simple (below the self-window)
        without displaying PAMPs. The immune ``blind spot'' for amyloid
        aggregates applies generally to protein aggregation diseases.

  \item \textbf{Rich vs.\ Depth Orthogonality (Proposition~\ref{prop:richness_vs_depth})}:
        $\Rich$ measures combinatorial complexity; $\Depth$ measures hierarchical
        autonomy. These are independent. Viruses are complex ($\Rich$ high)
        but shallow ($\Depth = 1$--$2$). Camouflage requires high $\Rich$;
        life requires high $\Depth$.

  \item \textbf{Innate as Anomaly Detector (Theorem~\ref{thm:innate_anomaly})}:
        Innate immunity detects categorical geometric anomalies (LPS architecture,
        dsRNA topology, cytoplasmic DNA location) in dimensions absent from
        self-space. Cross-reactivity is a feature, not a limitation.

  \item \textbf{Three-Layer Evasion Stack (Definition~\ref{def:evasion_stack})}:
        Successful viruses deploy protein richness camouflage, assembly
        minimisation, and antigenic drift simultaneously. Incubation period is
        the time to cross the assembly-geometry detection threshold.
\end{enumerate}
\end{chapterbox}
